\QHeading{第1题　口算。}
\par\noindent


\begin{oralgrid}
\OralItem{1}{0.75+0.25}{1}
\OralItem{2}{2.3+1.8}{4.1}
\OralItem{3}{7.5-1.5}{6}
\OralItem{4}{8.3-2.3}{6}
\OralItem{5}{3.1+0.8}{3.9}
\OralItem{6}{4+0.68}{4.68}
\OralItem{7}{4.25+0.5}{4.75}
\OralItem{8}{3.76-0.6}{3.16}
\OralItem{9}{0.26+0.35}{0.61}
\OralItem{10}{0.37+14.3}{14.67}
\OralItem{11}{10-2.3}{7.7}
\OralItem{12}{5.01-1.07}{3.94}\end{oralgrid}
\par

\QHeading{口算思路补充（对应上表题号）}

\Answer{上面列出得数，下面按小题给出不同计算思路。计数单位法适用于小数，整数题可比较凑整、拆分或商不变等方法。}
\MethodTitle{第1小题：$0.75+0.25$}

\MethodBlock{口算思路一：按相同计数单位计算}{先看成若干个百分之一，按整数加减，最后还原成小数。}{\begin{aligned}[t]&\phantom{={}}0.75+0.25\\
&=(75+25)\div 100\\
&=100\div 100\\
&=1\end{aligned}}
\MethodTitle{第2小题：$2.3+1.8$}

\MethodBlock{口算思路一：按相同计数单位计算}{先看成若干个十分之一，按整数加减，最后还原成小数。}{\begin{aligned}[t]&\phantom{={}}2.3+1.8\\
&=(23+18)\div 10\\
&=41\div 10\\
&=4.1\end{aligned}}

\MethodBlock{口算思路二：凑整再减去多加的数}{把第二个数凑成整数，再减去多加的部分。}{\begin{aligned}[t]&\phantom{={}}2.3+1.8\\
&=2.3+2-0.2\\
&=4.3-0.2\\
&=4.1\end{aligned}}
\MethodTitle{第3小题：$7.5-1.5$}

\MethodBlock{口算思路一：按相同计数单位计算}{先看成若干个十分之一，按整数加减，最后还原成小数。}{\begin{aligned}[t]&\phantom{={}}7.5-1.5\\
&=(75-15)\div 10\\
&=60\div 10\\
&=6\end{aligned}}

\MethodBlock{口算思路二：补差法}{从减数先补到一个整数，再补到被减数，两段差相加。}{\begin{aligned}[t]&\phantom{={}}7.5-1.5\\
&=(2-1.5)+(7.5-2)\\
&=0.5+5.5\\
&=6\end{aligned}}
\MethodTitle{第4小题：$8.3-2.3$}

\MethodBlock{口算思路一：按相同计数单位计算}{先看成若干个十分之一，按整数加减，最后还原成小数。}{\begin{aligned}[t]&\phantom{={}}8.3-2.3\\
&=(83-23)\div 10\\
&=60\div 10\\
&=6\end{aligned}}

\MethodBlock{口算思路二：补差法}{从减数先补到一个整数，再补到被减数，两段差相加。}{\begin{aligned}[t]&\phantom{={}}8.3-2.3\\
&=(3-2.3)+(8.3-3)\\
&=0.7+5.3\\
&=6\end{aligned}}
\MethodTitle{第5小题：$3.1+0.8$}

\MethodBlock{口算思路一：按相同计数单位计算}{先看成若干个十分之一，按整数加减，最后还原成小数。}{\begin{aligned}[t]&\phantom{={}}3.1+0.8\\
&=(31+8)\div 10\\
&=39\div 10\\
&=3.9\end{aligned}}

\MethodBlock{口算思路二：凑整再减去多加的数}{把第二个数凑成整数，再减去多加的部分。}{\begin{aligned}[t]&\phantom{={}}3.1+0.8\\
&=3.1+1-0.2\\
&=4.1-0.2\\
&=3.9\end{aligned}}
\MethodTitle{第6小题：$4+0.68$}

\MethodBlock{口算思路一：按相同计数单位计算}{先看成若干个百分之一，按整数加减，最后还原成小数。}{\begin{aligned}[t]&\phantom{={}}4+0.68\\
&=(400+68)\div 100\\
&=468\div 100\\
&=4.68\end{aligned}}
\MethodTitle{第7小题：$4.25+0.5$}

\MethodBlock{口算思路一：按相同计数单位计算}{先看成若干个百分之一，按整数加减，最后还原成小数。}{\begin{aligned}[t]&\phantom{={}}4.25+0.5\\
&=(425+50)\div 100\\
&=475\div 100\\
&=4.75\end{aligned}}
\MethodTitle{第8小题：$3.76-0.6$}

\MethodBlock{口算思路一：按相同计数单位计算}{先看成若干个百分之一，按整数加减，最后还原成小数。}{\begin{aligned}[t]&\phantom{={}}3.76-0.6\\
&=(376-60)\div 100\\
&=316\div 100\\
&=3.16\end{aligned}}

\MethodBlock{口算思路二：补差法}{从减数先补到一个整数，再补到被减数，两段差相加。}{\begin{aligned}[t]&\phantom{={}}3.76-0.6\\
&=(1-0.6)+(3.76-1)\\
&=0.4+2.76\\
&=3.16\end{aligned}}
\MethodTitle{第9小题：$0.26+0.35$}

\MethodBlock{口算思路一：按相同计数单位计算}{先看成若干个百分之一，按整数加减，最后还原成小数。}{\begin{aligned}[t]&\phantom{={}}0.26+0.35\\
&=(26+35)\div 100\\
&=61\div 100\\
&=0.61\end{aligned}}
\MethodTitle{第10小题：$0.37+14.3$}

\MethodBlock{口算思路一：按相同计数单位计算}{先看成若干个百分之一，按整数加减，最后还原成小数。}{\begin{aligned}[t]&\phantom{={}}0.37+14.3\\
&=(37+1430)\div 100\\
&=1467\div 100\\
&=14.67\end{aligned}}
\MethodTitle{第11小题：$10-2.3$}

\MethodBlock{口算思路一：按相同计数单位计算}{先看成若干个十分之一，按整数加减，最后还原成小数。}{\begin{aligned}[t]&\phantom{={}}10-2.3\\
&=(100-23)\div 10\\
&=77\div 10\\
&=7.7\end{aligned}}

\MethodBlock{口算思路二：补差法}{从减数先补到一个整数，再补到被减数，两段差相加。}{\begin{aligned}[t]&\phantom{={}}10-2.3\\
&=(3-2.3)+(10-3)\\
&=0.7+7\\
&=7.7\end{aligned}}
\MethodTitle{第12小题：$5.01-1.07$}

\MethodBlock{口算思路一：按相同计数单位计算}{先看成若干个百分之一，按整数加减，最后还原成小数。}{\begin{aligned}[t]&\phantom{={}}5.01-1.07\\
&=(501-107)\div 100\\
&=394\div 100\\
&=3.94\end{aligned}}

\MethodBlock{口算思路二：补差法}{从减数先补到一个整数，再补到被减数，两段差相加。}{\begin{aligned}[t]&\phantom{={}}5.01-1.07\\
&=(2-1.07)+(5.01-2)\\
&=0.93+3.01\\
&=3.94\end{aligned}}


\par\noindent

\QHeading{第2题　脱式计算，能简算的要简算。}

\Problem{（1）$7.46-1.48-3.52$}
\par


\MethodTitle{方法一}
\Answer{{\Formula{\begin{aligned}[t]&=7.46-(1.48+3.52)\\&=7.46-5\\&=2.46\end{aligned}}}}

\MethodBlock{方法二：按运算顺序直接计算（对照）}{不先用运算律变形，按括号、乘除、加减的顺序计算。题目要求简算时，优先使用上面的简便解法。}{\begin{aligned}[t]&\phantom{={}}7.46-1.48-3.52\\
&=5.98-3.52\\
&=2.46\end{aligned}}
\Problem{（2）$1.08+12.7+0.92$}

\MethodTitle{方法一}
\Answer{{\Formula{\begin{aligned}[t]&=1.08+0.92+12.7\\&=(1.08+0.92)+12.7\\&=2+12.7\\&=14.7\end{aligned}}}}

\MethodBlock{方法二：先处理括号或乘除，再凑整再减去多加的数}{把第二个数凑成整数，再减去多加的部分。}{\begin{aligned}[t]&\phantom{={}}1.08+12.7+0.92\\
&=13.78+0.92\\
&=13.78+1-0.08\\
&=14.78-0.08\\
&=14.7\end{aligned}}
\Problem{（3）$7.36-(3.5+2.36)$}

\MethodTitle{方法一}
\Answer{{\Formula{\begin{aligned}[t]&=7.36-(2.36+3.5)\\&=7.36-2.36-3.5\\&=5-3.5\\&=1.5\end{aligned}}}}

\MethodBlock{方法二：去括号后计算}{加号后的括号去掉不变号，减号后的括号去掉要改变各项的加减号。}{\begin{aligned}[t]&\phantom{={}}7.36-(3.5+2.36)\\
&=7.36-3.5-2.36\\
&=3.86-2.36\\
&=1.5\end{aligned}}

\MethodBlock{方法三：先处理括号或乘除，再多减再加回}{减去凑成的整数后，把多减的部分加回来。}{\begin{aligned}[t]&\phantom{={}}7.36-(3.5+2.36)\\
&=7.36-5.86\\
&=7.36-6+0.14\\
&=1.36+0.14\\
&=1.5\end{aligned}}
\Problem{（4）$3.24-(0.24+2.1)$}

\MethodTitle{方法一}
\Answer{{\Formula{\begin{aligned}[t]&=3.24-0.24-2.1\\&=3-2.1\\&=0.9\end{aligned}}}}

\MethodBlock{方法二：按运算顺序直接计算（对照）}{不先用运算律变形，按括号、乘除、加减的顺序计算。题目要求简算时，优先使用上面的简便解法。}{\begin{aligned}[t]&\phantom{={}}3.24-(0.24+2.1)\\
&=3.24-2.34\\
&=0.9\end{aligned}}
\Problem{（5）$79.37+8.64-7.37$}

\MethodTitle{方法一}
\Answer{{\Formula{\begin{aligned}[t]&=79.37-7.37+8.64\\&=(79.37-7.37)+8.64\\&=72+8.64\\&=80.64\end{aligned}}}}

\MethodBlock{方法二：按运算顺序直接计算（对照）}{不先用运算律变形，按括号、乘除、加减的顺序计算。题目要求简算时，优先使用上面的简便解法。}{\begin{aligned}[t]&\phantom{={}}79.37+8.64-7.37\\
&=88.01-7.37\\
&=80.64\end{aligned}}
\Problem{（6）$7.45+(0.78+5.55)$}

\MethodTitle{方法一}
\Answer{{\Formula{\begin{aligned}[t]&=7.45+(5.55+0.78)\\&=(7.45+5.55)+0.78\\&=13+0.78\\&=13.78\end{aligned}}}}

\MethodBlock{方法二：去括号后计算}{加号后的括号去掉不变号，减号后的括号去掉要改变各项的加减号。}{\begin{aligned}[t]&\phantom{={}}7.45+(0.78+5.55)\\
&=7.45+0.78+5.55\\
&=8.23+5.55\\
&=13.78\end{aligned}}

\MethodBlock{方法三：按运算顺序直接计算（对照）}{不先用运算律变形，按括号、乘除、加减的顺序计算。题目要求简算时，优先使用上面的简便解法。}{\begin{aligned}[t]&\phantom{={}}7.45+(0.78+5.55)\\
&=7.45+6.33\\
&=13.78\end{aligned}}
\par\noindent

\QHeading{第3题　解决问题。}

\Problem{① 一根彩带长20米，第一次用去5.35米，第二次用去6.8米，还剩多少米？}
\par


\MethodTitle{方法一}
\Answer{{\Formula{\begin{aligned}[t]&\phantom{={}}20-5.35-6.8\\&=20-(5.35+6.8)\\&=20-12.15\\&=7.85\text{（米）}\end{aligned}}}\par 答：还剩7.85米。}

\MethodTitle{方法二：按使用先后逐次求剩余}
\Answer{先减第一次用去的，再减第二次用去的。}
\Answer{\Formula{\begin{aligned}[t]&\phantom{={}}20-5.35-6.8\\
&=14.65-6.8\\
&=7.85\text{（米）}\end{aligned}}}
\Answer{答：还剩7.85米。}
\Problem{② 面粉厂第一季度生产了69.05万袋面粉，第二季度生产了78.35万袋面粉，第三季度生产的比前两个季度的总和少6.7万袋。第三季度生产了多少万袋面粉？}

\MethodTitle{方法一}
\Answer{{\Formula{\begin{aligned}[t]&\phantom{={}}69.05+78.35-6.7\\&=147.40-6.70\\&=140.7\text{（万袋）}\end{aligned}}}\par 答：第三季度生产140.7万袋面粉。}

\MethodTitle{方法二：先从第二季度中减去差额，再加第一季度}
\Answer{前两个季度的总和减6.7万袋，也可以先让第二季度的数量减6.7万袋，再加第一季度。}
\Answer{\Formula{\begin{aligned}[t]&\phantom{={}}69.05+(78.35-6.7)\\
&=69.05+71.65\\
&=140.7\text{（万袋）}\end{aligned}}}
\Answer{答：第三季度生产140.7万袋面粉。}
